\chapter{Introducción}

Se implementó un Pioneer P3DX en CoppeliaSim y Lua capaz de navegar autónomamente
en una celda robotizada con evitación de obstáculos, conforme a los requisitos
de la Actividad~2~\cite{viuActividad2}. La escena
representa un caso de fábrica sencillo: R1 asiste a B1, llamado Bill, una persona
tomada del modelo disponible en CoppeliaSim y usada como \textit{mannequin}
operativo. Los fundamentos de robótica móvil diferencial y navegación autónoma
sobre los que se construye el sistema están descritos en \cite{siegwart2011}. El robot lleva piezas o herramientas a dos mesas, espera mientras
Bill trabaja, las recoge, las devuelve a su estantería y termina en carga.

El ciclo básico incluye coordinación temporal, evitación local, diferenciación
entre estantería y bloqueo, retorno de herramienta y regreso a carga. La escena
principal se guardó como:

Las cuatro variantes de estimación y mapeo implementadas (\texttt{KALMAN\_LANDMARK},
\texttt{GMAPPING\_GRID}, \texttt{HECTOR\_GRID\_MATCHING} y
\texttt{CARTOGRAPHER\_SUBMAP}) son implementaciones propias en Lua, inspiradas
conceptualmente en los algoritmos del mismo nombre pero sin sus optimizadores
originales: la variante de GMapping implementa una grid de ocupación log-odds
sin filtro de partículas Rao-Blackwellizado; la de Hector realiza una búsqueda
discreta de pose en lugar de ajuste Gauss-Newton multi-resolución; la de Cartographer
aplica cierre de ciclo local en lugar de optimización global de grafo de poses.
La comparación evalúa cuatro \emph{representaciones de mapa} bajo las mismas
condiciones de misión, no la equivalencia con los paquetes ROS originales.

\begin{center}
\texttt{actividad2/coppeliasim/Sim\_T2\_Phase1\_Basic.ttt}
\end{center}

\section{Objetivos}

\begin{itemize}
  \item Implementar un Pioneer P3DX diferencial en CoppeliaSim y Lua que navegue de forma autónoma en una celda de almacén.
  \item Seguir al \textit{mannequin} Bill con control de distancia y rumbo en la escena de seguimiento, comparando los modos P, PI, PID, LQR (\textit{Linear Quadratic Regulator}) y NMPC (\textit{Nonlinear Model Predictive Control}) (escena \path{Actividad2_1_Basic_Follower.ttt}); la escena warehouse usa PID como controlador base.
  \item Planificar y ejecutar las entregas de T1 y T2 con una máquina de estados reproducible y validable.
  \item Evitar obstáculos en tiempo real con campos potenciales reactivos y recuperación ante bloqueos locales.
  \item Comparar cuatro variantes de estimación y mapeo simultáneo SLAM (\textit{Simultaneous Localisation and Mapping}) en Lua —Kalman-landmark, grid log-odds, ajuste por scan-matching y submapas— con métricas exportadas.
\end{itemize}

\section{Relación con el enunciado}

La guía pide programar un Pioneer P3DX en CoppeliaSim/Lua, primero para avanzar
hacia un \textit{mannequin} mediante campos potenciales y después para esquivar
obstáculos mientras mantiene ese comportamiento. La entrega también exige la
integración en una celda robotizada, un PDF y una presentación.

En este trabajo el \textit{mannequin} se implementa como Bill/B1, el modelo de
persona móvil disponible en CoppeliaSim. R1, el Pioneer, no solo se acerca a
Bill: sigue una referencia detrás de la persona, transporta T1 y T2, respeta
esperas de trabajo humano, evita pallet, pilar, caja y obstáculos de Phase 2, y
vuelve a la estación de carga. Se añade una capa de validación cuantitativa con métricas exportadas, comparación
de controladores y mapeo local (Tabla~\ref{tab:traceability-guide}).

\begin{table}[H]
\centering
\small
\begin{tabularx}{\textwidth}{@{}P{0.33\textwidth}Y@{}}
\toprule
\textbf{Punto de la guía} & \textbf{Dónde queda cubierto}\\
\midrule
Seguir al \textit{mannequin} & Bill/B1 actúa como objetivo móvil y como solicitante de herramientas en la celda.\\
Campos potenciales & Capítulo de campos potenciales: atracción hacia Bill y repulsión por sensores.\\
Anticolisión & Lectura de 16 sensores de proximidad, reducción de velocidad, giro reactivo y parada crítica.\\
Integración en celda & Warehouse con racks, mesas WS1/WS2, herramientas T1/T2, estación de carga y máquina de estados.\\
Ampliación del script base & Seguimiento P/PI/PID/LQR/NMPC, estimación Kalman, variantes SLAM y Phase 2 con mapa desconocido.\\
\bottomrule
\end{tabularx}
\caption{Trazabilidad entre el enunciado de la Actividad 2 y la solución implementada.}
\label{tab:traceability-guide}
\end{table}

\section{Alcance de los algoritmos}

Los algoritmos corren dentro de CoppeliaSim y Lua. Su alcance es acotado respecto
a paquetes ROS completos: Kalman-landmark combina odometría, corrección de pose
con mediciones ruidosas y rasgos compactos de obstáculo. Las variantes inspiradas en GMapping, Hector y Cartographer
implementan respectivamente: mapeo por grid log-odds, corrección de pose
por scan-matching discreto y gestión de submapas con cierre de ciclo local.
Ninguna replica el estimador probabilístico original (filtro de partículas,
Gauss-Newton o grafo de poses), pero mantiene la representación de mapa
característica de cada enfoque.

\section{Vista general del resultado}

La Fig.~\ref{fig:phase1-layout} muestra la vista superior de la celda base, con R1, B1, las estanterías de T1 y T2, las mesas de trabajo y los obstáculos presentes desde el inicio.

\begin{figure}[H]
  \centering
  \includegraphics[width=0.96\textwidth]{phase1_layout_diagram.png}
  \caption{Vista superior de la celda base \texttt{Sim\_T2\_Phase1\_Basic}.}
  \label{fig:phase1-layout}
\end{figure}
